\documentclass[11pt,letterpaper]{article}

\usepackage[margin=1in]{geometry}
\usepackage{amsmath,amssymb}

\setlength{\parindent}{0pt}
\setlength{\parskip}{0.5em}

\title{%
  Assignment 1: Poisson--Gamma Posterior\\[0.4em]
  {\large CSDS 611 --- Advanced Topics and Methods of Data Science I}%
}
\author{Andrei Cozma}
\date{}

\begin{document}

\maketitle

\section*{Model and parameterization}

Consider the Bayesian model
\[
  w_n\mid\lambda\sim\operatorname{Poisson}(\lambda),
  \qquad n=1,\ldots,N,
  \qquad
  \lambda\sim\operatorname{Gamma}(\phi,\psi).
\]
Assume that the observations are conditionally independent given $\lambda$,
with $w_n\in\{0,1,2,\ldots\}$ and $\phi,\psi>0$.
Gamma distribution is parameterized by
shape $\phi$ and scale $\psi$.
Let $\mathbf{w}=(w_1,\ldots,w_N)$ and $S=\sum_{n=1}^{N}w_n$.
We seek the posterior density $p(\lambda\mid\mathbf{w})$.

\section*{1.
  Likelihood}

The Poisson probability mass function is
\[
  p(w_n\mid\lambda)
  =\frac{\lambda^{w_n}e^{-\lambda}}{w_n!}.
\]
Conditional independence gives the joint likelihood
\begin{align*}
  p(\mathbf{w}\mid\lambda)
   & =\prod_{n=1}^{N}
  \frac{\lambda^{w_n}e^{-\lambda}}{w_n!} \\
   & =\frac{\lambda^{S}e^{-N\lambda}}
  {\prod_{n=1}^{N}w_n!}.
\end{align*}

\section*{2.
  Prior}

The Gamma prior has density
\[
  p(\lambda)
  =\frac{1}{\Gamma(\phi)\psi^{\phi}}
  \lambda^{\phi-1}e^{-\lambda/\psi},
  \qquad \lambda>0.
\]

\section*{3. Apply Bayes' rule}

By Bayes' rule,
\begin{align*}
  p(\lambda\mid\mathbf{w})
   & \propto p(\mathbf{w}\mid\lambda)\,p(\lambda) \\
   & =\frac{\lambda^{S}e^{-N\lambda}}
  {\prod_{n=1}^{N}w_n!}
  \frac{\lambda^{\phi-1}e^{-\lambda/\psi}}
  {\Gamma(\phi)\psi^{\phi}}                       \\
   & \propto\lambda^{\phi+S-1}
  \exp\!\left[-\left(N+\frac{1}{\psi}\right)\lambda\right].
\end{align*}
The final proportionality omits only factors that do not depend on $\lambda$.

\section*{4.
  Normalize and identify the posterior}

Introduce a normalizing constant $C$ and write
\[
  p(\lambda\mid\mathbf{w})
  =C\lambda^{\phi+S-1}
  \exp\!\left[-\left(N+\frac{1}{\psi}\right)\lambda\right],
  \qquad \lambda>0.
\]
The posterior must integrate to one. Using the substitution
$t=(N+1/\psi)\lambda$, we obtain
\begin{align*}
  1
   & =\int_{0}^{\infty}p(\lambda\mid\mathbf{w})\,d\lambda \\
   & =C\int_{0}^{\infty}\lambda^{\phi+S-1}
  e^{-(N+1/\psi)\lambda}\,d\lambda                        \\
   & =C\left(N+\frac{1}{\psi}\right)^{-(\phi+S)}
  \int_{0}^{\infty}t^{\phi+S-1}e^{-t}\,dt.
\end{align*}
The remaining integral equals $\Gamma(\phi+S)$, by normalization of the
Gamma density with shape $\phi+S$ and scale $1$.
Therefore,
\[
  1=C\frac{\Gamma(\phi+S)}
  {\left(N+\frac{1}{\psi}\right)^{\phi+S}},
  \qquad
  C=\frac{\left(N+\frac{1}{\psi}\right)^{\phi+S}}
  {\Gamma(\phi+S)}.
\]
Substituting $C$ gives the normalized posterior density:
\[
  \boxed{
    p(\lambda\mid\mathbf{w})
    =\frac{\left(N+\frac{1}{\psi}\right)^{\phi+S}}
    {\Gamma(\phi+S)}
    \lambda^{\phi+S-1}
    \exp\!\left[-\left(N+\frac{1}{\psi}\right)\lambda\right],
    \quad \lambda>0.
  }
\]
Comparing this expression with the Gamma density gives the updated shape
and scale:
\[
  \phi'=\phi+\sum_{n=1}^{N}w_n,
  \qquad
  \psi'=\frac{1}{N+1/\psi}
  =\frac{\psi}{1+N\psi}.
\]
Thus,
\[
  \boxed{
    \lambda\mid w_1,\ldots,w_N
    \sim\operatorname{Gamma}\!\left(
    \phi+\sum_{n=1}^{N}w_n,
    \frac{\psi}{1+N\psi}
    \right)
    \quad\text{(shape--scale).}
  }
\]
The posterior remains in the Gamma family, so the Gamma prior is conjugate
to the Poisson likelihood.

\end{document}
